\section{Ciphertext Integrity}
\label{sec:authenticity}

This section proves the random-oracle INT-CTXT bound feeding the
all-in-one AE composition of \Cref{thm:ae}. The public-permutation
INT-CTXT statement follows by the flat-sponge lift of
\Cref{app:flat-sponge-lift}.

\begin{theorem}[Random-Oracle INT-CTXT]
\label{thm:intctxt-ro}
For every INT-CTXT adversary $\mathcal{B}$ in the \TW{} random-oracle
model (\Cref{sec:ro-model}),
\[
  \Adv^{\mathsf{int\text{-}ctxt}}_\RO(\mathcal{B})
  \;\leq\;
  \KeyGuess_k(\qRO(\mathcal{B}))
  + \frac{\nleaf(\mathcal{B})\bigl(\nleaf(\mathcal{B}) - 1\bigr)}{2^{\tauleaf + 1}}
  + \frac{\qv(\mathcal{B})}{2^{\tauroot}}.
\]
Equivalently, the resource-indexed envelope satisfies
\[
  \Epsintctxt(\QRO, \Nleaf, \Qv)
  \;\leq\;
  \KeyGuess_k(\QRO)
  + \frac{\Nleaf(\Nleaf - 1)}{2^{\tauleaf + 1}}
  + \frac{\Qv}{2^{\tauroot}}.
\]
\end{theorem}

\begin{proof}
Let $\mathsf{G}_0$ be the real random-oracle INT-CTXT game: the
challenger samples $K \sample \bits^k$, $\mathcal{B}$ has direct
access to $\ROflat$, and construction queries are answered by
$\Enc_K^\RO$ and $\Ver_K^\RO$. Each encryption query
$(N, A, M)$ returns $C \concat T = \Enc_K^\RO(N, A, M)$ and records the
tuple $(N, A, C, T)$; $\mathcal{B}$ wins if any verification query
returns $1$ on a tuple $(N, A, C, T)$ that was not previously
recorded by the encryption oracle.

Let $\badkey$ be the event that some direct
adversary query $(Y,\ell)$ to $\ROflat$ satisfies
$\KeyedTWOut(Y) = (K,\mathcal{C})$ for a counted output-point class
$\mathcal{C}$ of \Cref{tab:keyed-twout}; the requested length $\ell$
is ignored. Let $\mathsf{G}_1$ be $\mathsf{G}_0$ with an immediate
abort on $\badkey$. The games are identical until
this event, so
\[
  \Pr[\mathsf{G}_0 \text{ wins}]
  \leq
  \Pr[\mathsf{G}_1 \text{ wins}]
  + \Pr[\badkey].
\]
\Cref{lem:int-key-uniformity} verifies the visible-prefix uniformity
invariant required by the Adaptive Key-Test Bound
(\Cref{lem:key-tests}) in $\mathsf{G}_1$, with the set of already
tested candidate keys recovered by $\KeyedTWOut$. Therefore
\[
  \Pr[\badkey]
  \;\leq\; \KeyGuess_k(\qRO(\mathcal{B})).
\]

Let $\badleaf$ be the event that two distinct
construction-generated final leaf transcripts under $K$ receive the
same $\tauleaf$-bit $\RF$ projection. Let $\mathsf{G}_2$ be
$\mathsf{G}_1$ with an additional immediate abort on
$\badleaf$. The games are identical until this
event, so
\[
  \Pr[\mathsf{G}_1 \text{ wins}]
  \leq
  \Pr[\mathsf{G}_2 \text{ wins}]
  + \Pr[\badleaf \text{ before }
        \badkey].
\]
Before $\badkey$, no direct query has pre-read the bridged input of a
hidden-key final leaf transcript before the construction first creates
that transcript: such a query would be decoded by $\KeyedTWOut$ to the
hidden key $K$ and its $\Ltag[j]$ class, independently of the requested
output length, and would itself trigger $\badkey$. This supplies the
event premise for the Leaf Tag Collision Bound
(\Cref{lem:leaf-collision}): thus
$\{\badleaf \text{ before }
\badkey\}$ is contained in
$\{\badleaf \text{ before }
\mathsf{hit}_{\mathsf{leaf}}\}$, where $\mathsf{hit}_{\mathsf{leaf}}$
is the event of \Cref{lem:leaf-collision} that a direct query hits the
bridged input of a hidden-key final leaf transcript before the
construction first creates it; hence
\[
  \Pr[\badleaf \text{ before }
      \badkey]
  \;\leq\;
  \frac{\nleaf(\mathcal{B})\bigl(\nleaf(\mathcal{B}) - 1\bigr)}{2^{\tauleaf + 1}}.
\]

\paragraph{Bound on $\Pr[\mathsf{G}_2 \text{ wins}]$.}
In $\mathsf{G}_2$, every accepting non-replay verification submission
belongs to a final root transcript class. If the class is
encryption-first---first generated by an encryption query returning
$(N_e,A_e,C_e,T_e)$---then, since $\badleaf$ has not occurred, the
final-root determinacy part of \Cref{lem:transcript-inj} gives
$(N,A,C) = (N_e,A_e,C_e)$ and the correct root tag is the value already
returned as $T_e$, so the submission is a replay when $T = T_e$ and
fails the tag comparison otherwise; it cannot be an accepting
non-replay. The remaining accepting non-replay submissions are
therefore verification-first. By
\Cref{lem:int-verification-first-root-guess}, their total probability
before $\badkey \cup
\badleaf$ is bounded by the root tag guessing
term:
\[
  \Pr[\mathsf{G}_2 \text{ wins}]
  \;\leq\; \qv(\mathcal{B}) / 2^{\tauroot}.
\]

\paragraph{Conclusion.}
Combining the three stopped-game bounds,
\[
  \Adv^{\mathsf{int\text{-}ctxt}}_\RO(\mathcal{B})
  = \Pr[\mathsf{G}_0 \text{ wins}]
  \leq
  \KeyGuess_k(\qRO(\mathcal{B}))
  + \frac{\nleaf(\mathcal{B})\bigl(\nleaf(\mathcal{B}) - 1\bigr)}{2^{\tauleaf + 1}}
  + \frac{\qv(\mathcal{B})}{2^{\tauroot}}.
\]
Taking the supremum over INT-CTXT adversaries with
$\qRO(\mathcal{B}) \leq \QRO$, $\nleaf(\mathcal{B}) \leq \Nleaf$, and
$\qv(\mathcal{B}) \leq \Qv$ gives the envelope statement.
\end{proof}

\paragraph{Lineage.}
The random-oracle decomposition of the bound into a key-guessing term
and a tag-guessing term ($q\,2^{-k}$ and $2^{-t}$ in their notation)
follows the authenticity bound of the ideal scheme
$\mathcal{RO}_{\mathrm{wrap}}[\rho]$ in the SpongeWrap analysis
\cite[Lemma~6]{BDPV11}, which carries no sponge-capacity term. We realize
the key-guessing term as the $\badkey$ abort and
the tag-guessing term as the residual root tag guessing bound on
$\Pr[\mathsf{G}_2 \text{ wins}]$, with a \TW{}-specific leaf tag
birthday abort $\badleaf$ inserted between them.
The sponge-capacity differentiating loss of
\cite[\S5, Theorem~1]{BDPV11} does not appear at this random-oracle
layer; it enters only through the flat-sponge lift of
\Cref{app:flat-sponge-lift}, as the \cite{BDPV08} indifferentiability
term written $\Lift_c$. The \TW{}-specific steps come from
Transcript Injectivity (\Cref{lem:transcript-inj}), the Leaf Tag
Collision Bound (\Cref{lem:leaf-collision}), which bounds the
$\badleaf$ abort by a birthday on the
$\tauleaf$-bit leaf tag space, and the INT-CTXT
auxiliary lemmas
(\Cref{lem:int-key-uniformity,lem:int-verification-first-root-guess}),
with the leaf-freshness and encryption-first-replay steps argued inline.
